\chapter{Planeamento de vários jogos e integração dos módulos}
\label{chap:horizonte}

O módulo de planeamento de vários jogos de PLANTIA integra as componentes num plano de vários jogos,
com objetivo de vitórias e limites partilhados de minutos. O procedimento
\path{scripts/plan_horizon.py}, descrito em
\path{docs/HORIZON_PLANNING.md}, reutiliza os modelos e planeadores dos
módulos de previsão de resultados, avaliação de jogadores, planeamento de rotações e planeamento de substituições sobre os dados do módulo de preparação de dados.

\section{Informação congelada e objetivo competitivo}

Uma única data de corte, posterior ao treino dos modelos, delimita as
observações utilizáveis para todos os jogos. O calendário fornece datas
e equipas, sem resultados futuros. As variáveis de descanso consideram
os jogos agendados, mas estatísticas das equipas, número de jogos disputados
e Elo permanecem congelados. Trata-se de uma aproximação de força constante
para jogos futuros, distinta da avaliação diária do módulo de previsão de resultados.

O plantel e as prioridades individuais são calculados uma única vez na
data de corte. As prioridades usam o contexto de equipa da casa mesmo
nos encontros fora, permitindo uma pontuação de alocação consistente,
mas não uma nova previsão individual por encontro. Assume-se estabilidade
do plantel dentro da temporada, admitindo indisponibilidades explicitamente
fornecidas por jogo. O calendário histórico não é versionado, pelo que a
sua disponibilidade antecipada é uma hipótese do cenário retrospetivo.

Se \(p_g\) é a probabilidade de vitória da equipa planeada no jogo \(g\),
o número esperado de vitórias é \(\sum_g p_g\). Para calcular a probabilidade
de atingir pelo menos \(K\) vitórias, o sistema assume resultados
independentes e calcula a convolução de Bernoulli. Sendo \(q_g(k)\) a
probabilidade de \(k\) vitórias após \(g\) jogos,
\[
q_g(k)=(1-p_g)q_{g-1}(k)+p_gq_{g-1}(k-1),\qquad q_0(0)=1,
\]
com probabilidades nulas fora dos índices válidos. Assim,
\(\Pr(W\geq K)=\sum_{k=K}^{G}q_G(k)\).
A esperança não exige independência, mas esta distribuição exige-a.
Lesões comuns, forma da equipa e erros correlacionados dos modelos não
são representados. O objetivo de vitórias é incerto, não uma restrição
cujo cumprimento o solver possa garantir.

\section{Seleção conjunta de rotações}

Para cada jogo, o módulo de planeamento de rotações gera candidatos com limites individuais de
36, 32 e 28 minutos. Mantêm-se as restantes restrições de rotação, posições,
descanso e disponibilidade. Cada candidato viável cobre 240 minutos-jogador;
os candidatos inviáveis são excluídos.

Se \(z_{gc}\) indica a seleção do candidato \(c\) para o jogo \(g\),
\(m_{igc}\) os seus minutos e \(J_{gc}\) a sua utilidade, um segundo
modelo CP-SAT impõe
\[
\sum_c z_{gc}=1,\qquad
\sum_{g,c}m_{igc}z_{gc}\leq M_i.
\]
Para cada par de jogos em dias consecutivos, a soma dos minutos de cada
jogador não pode ultrapassar o limite correspondente. Por omissão,
\(M_i=96\) no horizonte de três jogos e o limite para dias consecutivos
é 64 minutos. Estes limites são hipóteses configuráveis do planeamento.

A função objetivo é
\[
\max\sum_{g,c}\left[1+4p_g(1-p_g)\right]J_{gc}z_{gc}.
\]
O fator dá prioridade a encontros mais equilibrados. É uma preferência
heurística, não uma estimativa de retorno do esforço em vitórias.
A probabilidade de alcançar o objetivo é reportada separadamente: a
seleção de rotações não altera as probabilidades previstas pelo módulo de previsão de resultados.

A optimalidade refere-se apenas aos candidatos gerados e aos coeficientes
inteiros arredondados. Cada candidato pode ser apenas viável, devido ao
limite de dez segundos na sua resolução; a seleção conjunta tem limite
de 30 segundos. Um conjunto inviável de candidatos não prova a inexistência
de outra solução fora desse conjunto.

\section{Políticas de substituição e replaneamento}

Após selecionar as rotações, o módulo de planeamento de substituições calcula uma política MDP por
jogo. Como conserva os minutos de cada jogador em cada período, conserva
também os limites acumulados do horizonte. A recompensa do MDP é distinta
da utilidade da rotação e ambas são apresentadas separadamente. O exemplo
exportado segue contexto equilibrado, embora a política cubra os três
contextos modelados. A extensão genética do módulo de planeamento de substituições permite comparar
uma política alternativa para cada rotação selecionada, mas não substitui
automaticamente a política MDP utilizada pelo horizonte integrado.

A execução produz previsões por jogo, distribuição do número de vitórias,
minutos acumulados, seleção dos candidatos e, por encontro, entradas,
solução, política, minutos e horário. Verificações independentes controlam
os limites individuais e os pares em dias consecutivos. Os testes
documentados incluem a distribuição de vitórias, o corte temporal comum,
conflitos de orçamento de minutos e calendários inválidos.

Depois de cada jogo, o utilizador atualiza observações e disponibilidade,
avança a data de corte e executa novamente o horizonte restante. Alterações
nos dados preparados exigem reconstruir os artefactos de desempenho devido
à verificação dos hashes. O procedimento não acompanha jogos em direto
nem converte resultados simulados em observações reais. Os limites cobrem
apenas os jogos fornecidos, pelo que o calendário deve incluir todos os
encontros da equipa no horizonte.
